\section{Exponential records: comparing two magnifications}
\label{sec:exponential}

Suppose $\lambda>1$. The scalar limit $u$ is then locally H\"older.
We magnify its variation to a nonconstant function on a punctured
sphere. The approximating level continua force every transform with a different
puncture to be constant wherever this function varies. A second
magnification makes two such transforms vary on the same two disks,
with disjoint value ranges, yielding a contradiction.
Scalar approximation indices throughout belong to the subsequence
of \Cref{prop:scalar-limit}.

\subsection{Magnifying variation}

\begin{lemma}[H\"older magnification]
\label{lem:holder-magnification}
Let $f$ be a real-valued, locally H\"older continuous function on an open set $\Omega\subset Z$.
Suppose $O$ is open, $\overline O\subset\Omega$, and $O$ contains a
compact path on which $f$ is nonconstant. Fix $z_0$ with
$\dist(z_0,O)>0$. There are depths $k_i\to\infty$, centers $x_i\in O$,
positive numbers $\sigma_i$, and expansions $h_i=v_i^{-1}$ at
$(x_i,k_i)$ such that, after passage to a subsequence,
\[
 p_i=h_i z_0\longrightarrow b,
 \qquad
 \frac{\dist(x_i,Z\setminus O)}{r_{k_i}}\longrightarrow\infty.
\]
Every compact subset of $Z\setminus\{b\}$ eventually lies in
$h_i\Omega$, and
\begin{equation}
 F_i(z)=\frac{f(h_i^{-1}z)-f(x_i)}{\sigma_i}
 \label{eq:blowup}
\end{equation}
converges locally uniformly there to a nonconstant locally H\"older
continuous function $F$.
\end{lemma}

\begin{proof}
Write $\delta(x)=\dist(x,Z\setminus O)$ and fix $\eta>Q$.
A finite cover of $\overline O$ gives a uniform H\"older bound,
with some exponent $\alpha_0>0$, on a neighborhood of $\overline O$.
For all large $k$ put
\[
 A_k=\sup_{x\in O}\delta(x)^\eta
            \osc_{\overline B(x,r_k)}f.
\]
The balls lie in $\Omega$, and
\begin{equation}
 c a^{-Qk}\le A_k\le C a^{-\alpha_0 k}.
 \label{eq:magnification-bounds}
\end{equation}
For the lower bound, cover a compact variation path $P\subset O$
by at most $C r_k^{-Q}$ radius-$r_k$ balls centered on $P$, using
\Cref{lem:packing}. Their range intervals cover the interval $f(P)$,
so one ball has oscillation at least $c r_k^Q$; also
$\min_P\delta>0$. The upper bound is the H\"older estimate.

Fix $0<\alpha<\alpha_0$. Since $A_k a^{\alpha k}\to0$, there are
arbitrarily late indices attaining their maximum over the remaining
tail. At these indices,
\begin{equation}
 A_{k+j}\le a^{-\alpha j}A_k\qquad(j\ge0).
 \label{eq:magnification-tail}
\end{equation}
Choose $x_k\in O$ attaining at least half the supremum, and set
$\delta_k=\delta(x_k)$ and $\sigma_k=A_k/\delta_k^\eta$.
Boundedness of $f$ and \eqref{eq:magnification-bounds} give
$\delta_k^\eta\ge c r_k^Q$, hence
$\delta_k/r_k\ge c r_k^{Q/\eta-1}\to\infty$.
Expand at $(x_k,k)$ and take a subsequence with $p_k=h_kz_0\to b$.

For a compact $K\subset Z\setminus\{b\}$, \eqref{eq:puncture} gives
\begin{equation}
 h_k^{-1}K\subset\overline B(x_k,C_Kr_k)\subset O
 \label{eq:magnification-preimage}
\end{equation}
eventually. The expansion formula gives
$d(h_k^{-1}z,h_k^{-1}w)\le C_Kr_kd(z,w)$ on $K$.
On every fixed multiple of this source ball, $\delta$ is comparable
to $\delta_k$. Thus \eqref{eq:magnification-tail} implies
\begin{equation}
 \frac{|f(y)-f(y')|}{\sigma_k}\le C a^{-\alpha j}
 \quad\text{if }d(y,y')\le r_{k+j},\quad j\ge0.
 \label{eq:magnification-local-holder}
\end{equation}
For sufficiently small $d(z,w)$, choosing $j\ge0$ with
$a^{-j}\ge C_Kd(z,w)>a^{-j-1}$ proves a uniform local
$\alpha$-H\"older estimate for $F_k$ on $K$.

For boundedness, \Cref{lem:paths} joins $h_k^{-1}z$ to $x_k$
inside a fixed multiple of $B(x_k,r_k)$. Packing covers this path
by a bounded number of radius-$r_k$ balls centered on the path, each with oscillation
at most $C\sigma_k$. Their range intervals cover the path's image,
so $|F_k(z)|\le C_K$. Compact exhaustion and Arzel\`a--Ascoli now
give a locally uniform, locally H\"older limit.
Finally, $\osc_{h_k\overline B(x_k,r_k)}F_k\ge1/2$.
The lower puncture bound places all these image balls in one compact
set off $b$, so local uniform convergence makes the limit nonconstant.
\end{proof}

\subsection{A common mesh and escaping level continua}

Apply \Cref{lem:holder-magnification} to $u$ on $U$, taking a small
variation region $O$ and $z_0$ outside its closure. Call the limit
$W$, its puncture $b$, and retain $x_i,k_i,h_i,\sigma_i$.

\begin{lemma}[Diagonal approximants]
\label{lem:exponential-diagonal}
After thinning the expansions, there are increasing scalar indices
$N_j$ such that
\[
 f_j(z)=\frac{H_{N_j}(h_j^{-1}z)-u(x_j)}{\sigma_j},
 \qquad z\in h_j\Qhat,
\]
converges locally uniformly to $W$ on $Z\setminus\{b\}$.
For every fixed $g\in G$, compact $K\subset Z\setminus\{gb\}$,
and fixed $J\ge2$,
\begin{equation}
 \mathcal E_{N_j}^{(J)}(f_j\circ g^{-1};K)\longrightarrow0.
 \label{eq:exponential-translated-energy}
\end{equation}
All indicated enlarged balls eventually lie in the functions' domains.

If a compact path $P\subset Z\setminus\{b\}$ has endpoint values
of $W$ strictly straddling a compact positive-length interval $I$,
then, for all large $j$ and every $t\in I$, there is a continuum
$D_{j,t}\subset h_j\Qhat$ with
\[
 f_j=t\text{ on }D_{j,t},\qquad D_{j,t}\cap P\ne\varnothing,
 \qquad \sup_{t\in I}\dist(b,D_{j,t})\longrightarrow0.
\]
\end{lemma}

\begin{proof}
For $j\ge4$ put
$K_j=\{d(z,b)\ge1/j\}$ and $L_j=\{d(z,b)\ge1/(2j)\}$.
Choose the $j$th expansion so that $h_j^{-1}L_j\subset U$ and its
normalized $u$-transform differs from $W$ by at most $1/j$ on $L_j$.
\emph{Hold this expansion and its normalization fixed.} Then choose
$N_j>N_{j-1}$ so large that
\[
 \sup_{L_j}\frac{|H_{N_j}(h_j^{-1}z)-u(h_j^{-1}z)|}{\sigma_j}
 \le\frac1j,
 \qquad \mathcal E_{N_j}^{(j)}(f_j;K_j)\le\frac1j,
 \qquad (j+2)r_{N_j}<\frac1{2j}.
\]
The uniform approximation follows from \Cref{prop:scalar-limit}.
With $h_j$ and $\sigma_j$ fixed, \Cref{lem:local-oscillation} bounds
the displayed energy by $C(h_j,K_j,j)/(\sigma_j^2M(N_j))$, so the
energy condition follows by increasing $N_j$. The last condition
keeps all the enlarged balls in $L_j$. This proves local uniform
convergence and vanishing energy at every fixed enlargement off $b$.

For a fixed $g$, choose a compact neighborhood $K'$ of $g^{-1}K$
avoiding $b$. Pulling back each target ball and choosing a nearest
source net point gives, by bi-Lipschitz distortion and packing,
\[
 \mathcal E_n^{(J)}(F\circ g^{-1};K)
 \le C_g\mathcal E_n^{(J')}(F;K')
\]
for some fixed $J'\ge2$ and all large $n$. Each source center receives
only boundedly many target centers. Eventually $K'\subset K_j$ and
$J'\le j$, proving \eqref{eq:exponential-translated-energy}.

Uniform convergence makes the endpoint values of $f_j$ straddle $I$.
The source level $\tau_j(t)=u(x_j)+\sigma_jt$ therefore lies in $(0,1)$,
and \Cref{lem:level-continua} supplies
$D_{j,t}=h_jC_{N_j,\tau_j(t)}$, meeting $P$.
Its source continuum has diameter at least $d_1$, so contains a point
at distance at least $d_1/2$ from $x_j$. By \eqref{eq:puncture},
\[
 \dist(b,D_{j,t})\le d(b,p_j)+Cr_{k_j}\longrightarrow0
\]
uniformly for $t\in I$.
\end{proof}

\subsection{Comparison across punctures}

Let $\mathcal N(F)$ be the points of a continuous function's domain
at which it is not locally constant.

\Needspace{5\baselineskip}
\begin{lemma}[Transform comparison]
\label{lem:exponential-comparison}
If $c=gb\ne b$ and $V=W\circ g^{-1}$, then
\begin{equation}
 V(x)=V(b)\qquad
 \bigl(x\in\mathcal N(W)\setminus\{c\}\bigr).
 \label{eq:transform-level}
\end{equation}
More generally, for $W_s=W\circ s^{-1}$ and $p_s=sb$, if $p_s\ne p_t$,
\begin{equation}
 W_t(x)=W_t(p_s)\qquad
 \bigl(x\in\mathcal N(W_s)\setminus\{p_t\}\bigr).
 \label{eq:two-transform-level}
\end{equation}
\end{lemma}

\begin{proof}
Suppose $x\in\mathcal N(W)\setminus\{c\}$ and $V(x)\ne V(b)$.
Choose $\varepsilon>0$ and a coordinate disk about $x$, compactly
contained off $b,c$, on which $|V-V(b)|\ge4\varepsilon$.
It contains a compact variation path $P$ for $W$. Choose a compact
positive-length interval $I$ strictly between its endpoint values,
avoiding $W(c)$. Continuity gives $\epsilon_c>0$ such that
$\overline B(c,\epsilon_c)$ avoids $b$ and its $W$-range has positive
distance from $I$. Consequently, for all large $j$, the continua
$D_{j,t}$ from \Cref{lem:exponential-diagonal} avoid this closed ball.

Choose $\rho>0$ so that $\overline B(b,\rho)$ avoids $c,P$ and
$|V-V(b)|\le\varepsilon$ there. For $w_j=f_j\circ g^{-1}$,
local uniform convergence gives
\[
 |w_j-V(b)|\ge3\varepsilon\quad\text{on }P,
 \qquad |w_j-V(b)|\le2\varepsilon\quad\text{on }\overline B(b,\rho)
\]
for all large $j$. Each $D_{j,t}$ meets $P$ and enters $B(b,\rho)$.
By \Cref{lem:stop-continuum}, it has a subcontinuum joining a chosen
point of $P$ to the radius-$\rho$ sphere and contained in
\[
 K=\{z:d(z,b)\ge\rho,\ d(z,c)\ge\epsilon_c\}.
\]
This subcontinuum lies in the level $f_j=t$ and has $w_j$-oscillation
at least $\varepsilon$. A fixed compact neighborhood of $K$ avoiding
both punctures eventually lies in $h_jU\cap gh_jU$ and contains every
closed radius-$2r_{N_j}$ cell meeting $K$. Hence
\Cref{lem:cross-energy} and Cauchy--Schwarz give
\[
 \varepsilon|I|\le
 \bigl(\mathcal E_{N_j}^{(2)}(f_j;K)\bigr)^{1/2}
 \bigl(\mathcal E_{N_j}^{(2)}(w_j;K)\bigr)^{1/2}\longrightarrow0,
\]
a contradiction. All sets and $g$ are fixed in this limit.
Finally apply \eqref{eq:transform-level} at $s^{-1}x$ with group
element $s^{-1}t$ to obtain \eqref{eq:two-transform-level}.
\end{proof}

\subsection{The second magnification}

We first control repeated punctures. An accumulation direction of
vertices going to infinity is a subsequential boundary limit.

\begin{lemma}[Stabilizer directions]
\label{lem:stabilizer}
For every $b\in Z$, the vertices of $E=\Stab_G(b)$ have at most
two accumulation directions.
\end{lemma}

\begin{proof}
Let $\xi$ be a based ray to $b$. For $q\in E$, $|q|\le t$,
thinness of the ideal triangle with sides $[e,q],\xi,q\xi$ puts
$q\xi(n)$ uniformly near $\xi$ when $n$ is sufficiently larger
than $t$. The nearby parameter lies in $[n-t-C,n+t+C]$.
For one common integer $n$, these vertices are distinct as $q$ varies.
Bounded valence therefore gives
\begin{equation}
 \#\{q\in E:|q|\le t\}\le C'(t+1).
 \label{eq:stabilizer-growth}
\end{equation}

If another accumulation direction exists, take $q_n\to x\ne b$ and
$q_n^{-1}\to y$ along a subsequence. Base change gives, uniformly
on compact sets avoiding $y$,
\[
 (q_n|q_nz)_e=|q_n|-(q_n^{-1}|z)_e+O(1)\ge |q_n|-O(1).
\]
By \Cref{lem:prefix}, $q_nz\to x$ there, and $q_nb=b$ forces $y=b$.
Choose disjoint small neighborhoods of $b,x$, with a point outside
their union. For large $n$, $q=q_n$ sends the complement of the
first into the second. Its positive iterates never return the outside
point, so $q$ has infinite order. It is therefore loxodromic, with
endpoints $b,x'$; see \cite[III.$\Gamma$.3.10]{BH}.

If some $h\in E$ moved $x'$, the cyclic subgroups generated by $q$
and $hqh^{-1}$ would intersect trivially: a nonzero common power
would have the same pair of endpoints. Their $m^2$ products
$q^i(hqh^{-1})^j$, $1\le i,j\le m$, would then be distinct and have
length $O(m)$, contradicting \eqref{eq:stabilizer-growth}.
Thus every element of $E$ fixes both endpoints. It maps a geodesic
line between them to a line at uniformly bounded Hausdorff distance,
so its orbit of $e$ stays a fixed distance from that line. Its only
possible accumulation directions are the two endpoints.
\end{proof}

\begin{proposition}
\label{prop:no-exponential-records}
The contradiction setup cannot have $\lambda>1$.
\end{proposition}

\begin{proof}
Let $E=\Stab(b)$ and let $\Lambda_E$ be its at-most-two accumulation
directions. Choose a variation region $O$ for $W$ whose closure avoids
$\{b\}\cup\Lambda_E$. Such a region exists: otherwise $W$ would
be locally constant on the connected sphere minus this finite set,
and continuity would make it constant throughout its domain.

Apply \Cref{lem:holder-magnification} to $W$ on $Z\setminus\{b\}$,
with this $O$ and $z_0=b$. Write the expansions as $s_i=v_i^{-1}$,
with centers $y_i\in O$, depths $\ell_i\to\infty$, and punctures
$p_i=s_i b\to b_*$. Positive affine normalizations of
$W_i=W\circ s_i^{-1}$ converge locally uniformly off $b_*$ to a
continuous nonconstant $W_*$.

The punctures have a pairwise distinct subsequence. Indeed, if
$v_i^{-1}b=v^{-1}b$ infinitely often, then $v_i v^{-1}\in E$.
After a subsequence $y_i\to y\in\overline O$. Since $v_i$ lies at
depth $\ell_i$ on a ray to $y_i$, \Cref{lem:prefix} gives $v_i\to y$
in boundary direction. Bounded right multiplication preserves this
direction, implying $y\in\Lambda_E$, a contradiction.

Choose two coordinate disks $D_1,D_2$, with compact closures off $b_*$,
having disjoint $W_*$-range intervals and each containing a compact
variation path. To do this, follow a variation path to its first
hitting points of two distinct intermediate values. Just before either
first hitting time the values differ, so the function is not locally
constant at either point. Continuity supplies disks with disjoint
ranges, and non-local constancy supplies a variation path in each.

For every sufficiently large $i$, the puncture $p_i$ is outside the
closed disks. Uniform convergence of the normalized $W_i$ preserves
both the disjoint range intervals and the positive variation on the
two fixed paths. Positive affine normalization preserves these
properties, so $\mathcal N(W_i)$ meets each disk.
Fix two distinct such indices $i,j$. By
\eqref{eq:two-transform-level}, $W_j$ has the single value $W_j(p_i)$
at every point of $\mathcal N(W_i)$ in either disk. Taking one point
in each contradicts the disjoint $W_j$-range intervals.
\end{proof}
